\documentclass[10pt,a4paper]{amsart}
\usepackage[margin=19mm]{geometry}
\usepackage{amsmath,amssymb,mathtools}
\usepackage[scr=rsfs,cal=euler]{mathalfa}
\usepackage{xfrac}
\usepackage[unicode,hidelinks,bookmarks=false]{hyperref}
\newcommand{\ie}{i.e.\ }
\newcommand{\cf}{cf.\ }
\newcommand{\resp}{respectively\ }
\newcommand{\Subsection}{\subsection}
\providecommand{\nobreakdash}{\nobreak\hbox{-}}
\setlength{\emergencystretch}{2em}
\multlinegap0pt
\usepackage{tikz}
\usepackage{mathtools}
\usepackage[shortlabels]{enumitem}
\usepackage{amssymb}
\usepackage{algorithm}\usepackage{algpseudocode}
\usepackage{xcolor}
\usepackage{dsfont}
\newtheorem{lemma}{Lemma}
\newcommand{\rev}[1]{{\color{black}#1}}
\title{Superadditivity-based valid inequalities and asymptotic bounds for the vehicle routing problem with stochastic demands}
\author{Robin Legault, Panca Jodiawan, Jean-Fran\c{c}ois C\^ot\'e, Leandro C. Coelho}
\date{Source: arXiv:2508.05877v2 (CC BY 4.0)}
\begin{document}
\maketitle

\noindent\emph{Source and licence.} Superadditivity-based valid inequalities and asymptotic bounds for the vehicle routing problem with stochastic demands. Version arXiv:2508.05877v2 (CC BY 4.0), distributed by arXiv under the Creative Commons Attribution 4.0 International licence, http://creativecommons.org/licenses/by/4.0/, as recorded in the arXiv licence metadata for this exact version and shown on the abstract page https://arxiv.org/abs/2508.05877v2. Full licence notice: \texttt{licenses/arxiv-2508.05877-CC-BY-4.0.txt}.\par\smallskip

\noindent\textit{Editorial scope.} Counted unit: the article's lemma lemma:monotonicity\_cost\_to\_go (source0.tex lines 505--508). The article's complete original proof of it is delivered with it (source0.tex lines 509--591, one \texttt{proof} environment of 83 lines). No mathematics is written, repaired or re-derived: the statement and the proof are the article's own lines, in the article's own notation, and no retained line is substituted or re-flowed.  Retained uncounted as the source-local context the proof needs: lines 266--274.  Nothing was omitted from any retained span, so the package carries no omission record.  No retained line contains a citation, so the wrapper carries no bibliography and no bibliography entry.  No retained line contains a figure environment, an image-inclusion command or drawing code, so no image is delivered and the package declares no asset dependency.  Editorial formatting only, in addition: an AMS \texttt{amsart} preamble that restates the article's own declarations and macros used by the retained lines, namely the environments \texttt{lemma} on separate counters, together with the standard packages those lines need.  The retained environments print in this document as Lemma 1 in order of appearance, whereas the article prints the same items with the numbering of its own full document.  The complete native source of the article as distributed in the arXiv e-print for this exact version is bundled unchanged as \texttt{sources/182902/source0.tex}.
\rev{Under the OR policy, after serving a customer, the vehicle may either perform a preventive restocking trip or proceed directly to the next customer. Here, for exact evaluation of the expected recourse cost by finite-state dynamic programming, we assume that the demands and vehicle capacity are integer-valued.} Denote by $\Psi(s,q) := \max\{\lceil(s - q)/Q\rceil,0\}$ the number of restocking trips required to fulfill a demand of~$s$ given residual capacity~$q$. For a route $(0, p, 0) = (i_0, i_{1}, \dots, i_{t}, i_{t+1})$, the expected recourse cost-to-go $F^{p}_{i_{j}}(q)$ when the vehicle leaves node $i_{j}$, for $j \in \{0,\dots,t{-}1\}$, with residual capacity $q \in \{0,\dots,Q\}$, is given by the following Bellman equation:
\begin{equation}
    F^p_{i_j}(q) := \min
    \begin{cases} \label{def:cost-to-go}
        H^p_{i_{j+1}}(q) := \sum\limits_{s \in \Xi_{i_{j+1}}} \left[c_{i_{j+1}}^F\Psi(s,q) + F^p_{i_{j+1}}\left(\Psi(s,q)Q + q -s  \right) \right]\rho_{i_{j+1}}^s, \\
        H^{*p}_{i_j,i_{j+1}} := c_{i_j, i_{j+1}}^P+H^p_{i_{j+1}}(Q),
    \end{cases}
\end{equation}
with the boundary condition $F^p_{i_{t}}(q):=0 \ \forall q \in \{0,\dots,Q\}$. This boundary reflects that, after serving the last customer $i_{t}$, the vehicle travels back to the depot as planned in the a priori solution, which does not incur a recourse cost. The terms $H^p_{i_{j+1}}(q)$ and $H^{*p}_{i_j,i_{j+1}}$ respectively denote the expected cost-to-go of proceeding directly to customer $i_{j+1}$ and of making a preventive return beforehand. Since travel distances respect the triangle inequality, proceeding to the next customer is always preferable when the vehicle has full capacity, i.e., $F^p_{i_j}(Q) = H^p_{i_{j+1}}(Q)$. We denote by $\bar{\mathcal{Q}}^{\text{OR}}_p:=F^p_{i_0}(Q)$ the expected recourse cost of route $(0,p,0)$ under the OR policy.

\begin{lemma}
\label{lemma:monotonicity_cost_to_go}
For any \rev{route $(0, p, 0) = (i_0, i_{1}, \dots, i_{t}, i_{t+1})$ and every position $j \in \{0,\dots,t\}$}, the expected recourse cost-to-go $F^{p}_{i_{j}}(q)$ is non-increasing in the residual capacity \rev{$q \in \{0,\dots,Q\}$}.
\end{lemma}

\begin{proof}
\rev{The proof is by reverse induction on $j\in\{0,\dots,t\}$. We first prove the claim directly for $j=t$ and $j=t-1$, and then show that, for each $j\in\{1,\dots,t-1\}$, monotonicity of $F^p_{i_j}(\cdot)$ implies monotonicity of
$F^p_{i_{j-1}}(\cdot)$.

The claim holds trivially for
$j=t$, since $F^p_{i_t}(q)=0$ for all $q \in \{0,\dots,Q\}$. We next verify the claim for
$j=t-1$. Since $F^p_{i_{t-1}}(q)
=
\min\{H^p_{i_t}(q),H^{*p}_{i_{t-1},i_t}\}$,
and the preventive-restocking term $H^{*p}_{i_{t-1},i_t}$ is independent of $q$, it is enough to
show that $H^p_{i_t}(\cdot)$ is non-increasing. Take integers $q_1,q_2$ with $0\le q_1<q_2\le Q$. Using that $F^p_{i_t}(q)=0$ for all $q \in \{0,\dots,Q\}$, we can write:
\[
H^p_{i_t}(q_1)-H^p_{i_t}(q_2)
=
\sum_{s\in\Xi_{i_t}}
c^F_{i_t}\bigl(\Psi(s,q_1)-\Psi(s,q_2)\bigr)\rho_{i_t}^s,
\]
which is non-negative since $\Psi(s,q)=\max\{\lceil(s - q)/Q\rceil,0\}$ is non-increasing in $q$ for any given $s$. Hence
$F^p_{i_{t-1}}(\cdot)$ is non-increasing. 

Now fix $j\in\{1,\dots,t-1\}$ and suppose that $F^p_{i_j}(\cdot)$ is non-increasing. We show that $F^p_{i_{j-1}}(\cdot)$ is non-increasing. Again, as $F^p_{i_{j-1}}(q) = \min\{H^p_{i_j}(q),H^{*p}_{i_{j-1},i_j}\}$,
and the preventive-restocking term is independent of $q$, it is enough to show that $H^p_{i_j}(\cdot)$ is non-increasing.

Take integers $q_1,q_2$ with $0\le q_1<q_2\le Q$. We show that $H^p_{i_j}(q_1)\ge H^p_{i_j}(q_2)$. For $\ell\in\{1,2\}$, denote by $\psi_\ell(s):=\Psi(s,q_\ell)$ and $r_\ell(s):=\psi_\ell(s)Q+q_\ell-s$ the number of restocking trips required to satisfy demand $s$ and the resulting residual capacity after serving customer $i_j$, respectively, when the vehicle arrives at $i_j$ with residual capacity $q_\ell$. Since $0\le q_1<q_2\le Q$, we have $\psi_1(s)-\psi_2(s)\in\{0,1\}$ for all $s\ge 0$. Thus, expanding the definition of $H^p_{i_j}$ and splitting the demand realizations according to whether starting from $q_2$ instead of $q_1$ reduces the number of required restocking trips, we obtain:
\begin{align}
H^p_{i_j}(q_1)-H^p_{i_j}(q_2)
=&
\sum_{s\in\Xi_{i_j}}
\left[
c^F_{i_j}(\psi_1(s)-\psi_2(s))
+
F^p_{i_j}(r_1(s))-F^p_{i_j}(r_2(s))
\right]\rho_{i_j}^s
\notag\\
=&
\sum_{\substack{s\in\Xi_{i_j}\\ \psi_1(s)=\psi_2(s)}}
\left[
F^p_{i_j}(r_1(s))-F^p_{i_j}(r_2(s))
\right]\rho_{i_j}^s
\notag\\
&+
\sum_{\substack{s\in\Xi_{i_j}\\ \psi_1(s)=\psi_2(s)+1}}
\left[
c^F_{i_j}
+
F^p_{i_j}(r_1(s))-F^p_{i_j}(r_2(s))
\right]\rho_{i_j}^s .
\label{eq:split}
\end{align}
The first sum in \eqref{eq:split} is non-negative by the induction hypothesis: when
$\psi_1(s)=\psi_2(s)$, we have $r_1(s)<r_2(s)$, and
$F^p_{i_j}(\cdot)$ is non-increasing.

It remains to bound the terms in the second sum. If
$\psi_1(s)=\psi_2(s)+1$, then $r_1(s)\in\{0,\dots,Q\}$, so the induction hypothesis
gives $F^p_{i_j}(r_1(s))\ge F^p_{i_j}(Q)$.
Moreover, by the Bellman recursion \eqref{def:cost-to-go} at node $i_j$, $F^p_{i_j}(r_2(s))
\le
H^{*p}_{i_j,i_{j+1}}
=
c^P_{i_j,i_{j+1}}+H^p_{i_{j+1}}(Q)$.
Therefore, each term in the second summation satisfies:
\begin{align*}
    c^F_{i_j} + F^p_{i_j}(r_1(s))-F^p_{i_j}(r_2(s)) & \geq c^F_{i_j}
+
F^p_{i_j}(Q)
-
c^P_{i_j,i_{j+1}}
-
H^p_{i_{j+1}}(Q) \\
&= c^F_{i_j}-c^P_{i_j,i_{j+1}} \\
&= \left(c_{0,i_j}+c_{i_j,i_{j+1}}-c_{0,i_{j+1}}\right)
+
\left(b^F-b^P\right) \\
& \geq 0.
\end{align*}
The first equality uses that $F^p_{i_j}(Q)=\min\{H^p_{i_{j+1}}(Q),c^P_{i_j, i_{j+1}} + H^p_{i_{j+1}}(Q)\}=H^p_{i_{j+1}}(Q)$, i.e., proceeding directly is always optimal at full capacity. The second equality uses the definitions of $c^F_{i_j}$ and $c^P_{i_j,i_{j+1}}$. The last inequality follows from the triangle inequality and
$b^F\ge b^P$. 

Hence both sums in \eqref{eq:split} are non-negative, and
$H^p_{i_j}(q_1)\ge H^p_{i_j}(q_2)$. We conclude that $H_{i_j}^p(\cdot)$ is non-increasing, and hence so is $F^p_{i_{j-1}}(\cdot)$.
}
\end{proof}

\end{document}
